\documentclass[11pt]{article}
\usepackage{amsmath,amssymb,booktabs}
\usepackage[margin=1in]{geometry}
\begin{document}

\section*{Step 212: Connes--Consani 2020 Audit for Branch A}

The audited source is Connes--Consani, \emph{Weil positivity and Trace formula
the archimedean place}, arXiv:2006.13771, using the manager-extracted text at
\texttt{/tmp/burnol\_audit/connes\_weil\_archimedean.txt}.

\subsection*{Translation}

\[
\begin{array}{lll}
\text{Cascade} & \text{Connes--Consani} & \text{Status}\\
\hline
K_\lambda \text{ at }\lambda=1 & S(1,1) & \text{identified}\\
P_\infty & S,\text{ projection onto }S(1,1) & \text{identified at }\lambda=1\\
M_{m_\ell} & \vartheta(f) & \text{related scaling action}\\
(I-P_\infty)M_{m_\ell}P_\infty & (I-S)\vartheta(f)S & \text{off-diagonal match}\\
q_\eta(C_\ell P_\eta) & \operatorname{Tr}(\vartheta(f)S), K_I,T_q & \text{not identified}
\end{array}
\]

\subsection*{Theorem 4.7}

Connes--Consani Theorem 4.7 concerns the positivity of
\[
  \operatorname{Tr}(\vartheta(f)S)
  =W_\infty(f)+\int f(\rho^{-1})\epsilon(\rho)\,d^*\rho.
\]
This is a compressed trace/positivity statement on the Sonin projection.
Branch A asks for the Calkin class of the off-diagonal object
\[
  q_\eta((I-S)\vartheta(f)S\,P_\eta).
\]
The trace of the compressed scaling action and the essential class of the
off-diagonal compression are different invariants.  Positivity of the former
does not decide compactness of the latter.

\subsection*{Section 6.2 Toeplitz approximation}

Section 6.2 discretizes a compact operator \(K_I\) by Toeplitz matrices \(T_q\).
It uses finite hermitian Toeplitz theory: largest eigenvectors, roots on the
unit circle, and a co-rank-one decomposition
\[
  S=\lambda\operatorname{Id} - \cdots T_q,
  \qquad
  S=\lambda\sum_j d(j)e(z_j).
\]
This is a finite approximation/canonicalization of \(K_I\).  It is not a
Toeplitz extension theorem for the infinite C*-algebra
\[
  A_\eta=C^*(P_\infty,M_{m_\ell},P_\eta,I),
\]
and does not give a faithful symbol
\(\sigma_B:A_\eta/K_\eta\to C(X)\).

\subsection*{Relation to Connes CRCFT-TE}

The paper remains valuable for the Connes/Weil positivity route.  In the
framework taxonomy it supports the Step 184 classification: the Connes/Weil
route is target-equivalent to RH, not a direct Branch A Calkin-symbol closure.

\subsection*{Verdict}

\[
\boxed{\texttt{V\_connes\_consani\_related\_but\_not\_sufficient}.}
\]

The Branch A G2 missing content remains a faithful boundary-symbol theorem for
the full \((P_\infty,M_{m_\ell},P_\eta)\) algebra.

\end{document}
